\section{Contratos OpenAPI y trazabilidad de interfaz}
\label{sec:contratos-openapi}

La E5 incorpora contratos OpenAPI 3.1.0 versionados para describir de forma
reproducible la superficie HTTP de los cuatro servicios de dominio y de la
fachada expuesta por el API Gateway. Los contratos se conservan en
\path{docs/openapi/}:

\begin{itemize}
  \item \texttt{auth-service-openapi.json};
  \item \texttt{usuarios-service-openapi.json};
  \item \texttt{academico-laboratorios-service-openapi.json};
  \item \texttt{reservas-solicitudes-service-openapi.json}; y
  \item \texttt{api-gateway-openapi.json}.
\end{itemize}

\subsection{Cobertura y fachada}

La fuente autoritativa contrasta cada snapshot con los
\texttt{RequestMappingInfo} que Spring registra en el
\texttt{RequestMappingHandlerMapping} del contexto de prueba. Auth cubre 12/12
operaciones, Usuarios 44/44 y Académico 71/71. Reservas cubre 80/80 con el
componente experimental ARBITER habilitado e incluye
\texttt{POST /api/v1/observabilidad/mobile/http-latency}; con su configuración
predeterminada hay 78 operaciones y no existen los dos mappings experimentales.
El contrato de 80 describe, por tanto, la superficie habilitable y no afirma que
ARBITER esté activo por defecto.

La fachada del Gateway contiene 225/225 operaciones método+ruta: 183 canónicas,
incluida \texttt{POST /api/v1/observabilidad/mobile/http-latency}, 8 alias de
Auth y 34 alias de Usuarios. El archivo \texttt{gateway-route-catalog.json}
actúa como lista blanca ejecutable: el Gateway solo acepta una petición si
\texttt{GatewayRouteCatalog.accepts(método, ruta)} reconoce ese par
método+ruta en el catálogo. Los seis beans de
\texttt{GatewayRoutes} delegan la construcción completa en
\texttt{rutaContratada}, que aplica obligatoriamente la conjunción del catálogo
con el alcance del servicio. El alcance solo recibe el path y no puede alterar
el resultado del catálogo. Se conservan los destinos y la reescritura de alias.
\texttt{RouterFunctionCatalogGuardTest} inspecciona bytecode con ASM: exige
la delegación en ese helper, rechaza construcciones o composiciones de routers
Spring fuera de él, incluso mediante llamadas indirectas, y verifica que el
helper alcanza \texttt{GatewayRouteCatalog.accepts}. No analiza operadores
booleanos ni deduce su semántica por mera alcanzabilidad.
\texttt{ContractedRouteTest} comprueba el router real: una ruta
\path{/api/v1/sin-contrato/recurso} es rechazada aunque el alcance devuelva
verdadero; una operación catalogada es aceptada con alcance compatible y
rechazada con alcance falso. La prueba de completitud comprueba los destinos
y las reescrituras para las operaciones documentadas.

\subsection{Seguridad contractual}

Las operaciones externas protegidas declaran un esquema HTTP Bearer con formato
JWT. Los endpoints públicos de autenticación no exigen ese esquema. Las APIs
internas de los servicios usan un esquema \texttt{apiKey} en el encabezado
\texttt{X-Internal-Api-Key}; esas rutas internas no forman parte del contrato de
fachada del Gateway.

\subsection{Producción, validación y CI}

Los snapshots se generan determinísticamente de controladores, DTO y rutas
vigentes sin requerir que los servicios arranquen con infraestructura externa.
Desde la raíz se validan o regeneran con:

\begin{lstlisting}
python scripts/validar-contratos-openapi.py
python scripts/validar-contratos-openapi.py --generate
\end{lstlisting}

El script Python conserva la generación y comprueba JSON, estructura OpenAPI,
identificadores de operación, referencias locales, esquemas de seguridad y
deriva del snapshot respecto de su productor. El mecanismo histórico obtenía
tanto el esperado como el snapshot mediante el mismo extractor regex; por ello
era circular como prueba de completitud. Ya no se presenta con ese alcance.

La completitud método+ruta se comprueba ahora de forma independiente con clases
compiladas y contextos Spring reales de prueba: no se leen fuentes Java como
texto ni se reutiliza \texttt{discover\_operations()}. El comparador es
bidireccional, normaliza solo los nombres de variables de path y rechaza rutas
faltantes o adicionales, método/path incorrectos y duplicados exactos o
semánticos. Pruebas negativas cubren además prefijos de clase,
\texttt{RequestMapping(method=...)}, mappings múltiples y deriva del Gateway.
El job matricial existente ejecuta \texttt{mvn verify} para los cinco módulos en
cada push y pull request; no existe un flag opcional que omita estas pruebas.

OpenAPI y Pact son complementarios. Pact verifica interacciones
consumidor--proveedor seleccionadas y no sustituye a OpenAPI. La validación
runtime acredita la completitud estructural método+ruta, pero no demuestra por
sí sola todos los comportamientos dinámicos ni las reglas de negocio durante
ejecución.

\subsection{Paginación, parámetros de consulta y validación independiente de schemas}

En los servicios, 25 operaciones GET devuelven resultados paginados: 14 usan
\texttt{Page<T>} de Spring Data (Académico y Usuarios) y 11 usan el DTO propio
\texttt{PaginaResponse<T>} (Reservas/Solicitudes/Agenda/Incidentes). El
extractor ya no describe estas respuestas como \texttt{type: array}; genera un
schema de objeto reutilizable por tipo (\texttt{Page\{Tipo\}} con
\texttt{content} y los metadatos que Spring serializa por defecto --
\texttt{totalElements}, \texttt{totalPages}, \texttt{number}, \texttt{size},
\texttt{numberOfElements}, \texttt{first}, \texttt{last}, \texttt{empty} --, o
\texttt{PaginaResponse\{Tipo\}} con los siete campos reales del \textit{record}
\texttt{PaginaResponse}: \texttt{contenido}, \texttt{pagina}, \texttt{tamanio},
\texttt{totalElementos}, \texttt{totalPaginas}, \texttt{primera} y
\texttt{ultima}).

El extractor tenía además un defecto que interpretaba el atributo
\texttt{defaultValue} de \texttt{@RequestParam} como si fuera el nombre del
parámetro, publicando parámetros de consulta llamados literalmente
\texttt{"0"}, \texttt{"20"} o \texttt{"60"} y marcados como obligatorios. Tras
corregir el parseo de atributos de anotación, los 23 parámetros afectados se
publican con su nombre real (\texttt{pagina}, \texttt{tamanio},
\texttt{rangoMinutos}), \texttt{required: false} y su \texttt{default} real. Los
endpoints que reciben \texttt{Pageable} documentan ahora \texttt{page},
\texttt{size} y \texttt{sort} como parámetros de consulta opcionales, acordes
con el comportamiento real de Spring Data Web: \texttt{page} y \texttt{size}
son enteros con defaults 0 y 20; \texttt{sort} es un array de strings sin
default. No hay overrides en las anotaciones ni en la configuración actual.

Esta corrección no puede acreditarse comparando el contrato contra sí mismo:
el extractor regex y el validador estructural comparten la misma
interpretación del código fuente, por lo que un contrato correcto podría ser
rechazado por un snapshot incorrecto generado por el mismo mecanismo, o un
defecto real podría pasar inadvertido si ambos lados heredan el mismo error.
Por eso la corrección de las respuestas paginadas, los parámetros de consulta
y los schemas de DTO se acredita con \texttt{SchemaContractVerifier}, una
clase Java independiente que usa reflexión sobre los \texttt{HandlerMethod}
que Spring ya resolvió. No lee las fuentes Java como texto ni reutiliza el
extractor Python como fuente de verdad: obtiene cuerpos de request y response,
\texttt{ResponseEntity}, \texttt{Page<T>} y \texttt{PaginaResponse<T>} desde
las clases compiladas mediante \texttt{ResolvableType}, y lee
\texttt{@RequestParam} o \texttt{Pageable} con \texttt{MethodParameter} para
obtener nombre, \texttt{required} efectivo, defaults y tipos simples. Las
pruebas \texttt{OpenApiSchemaContractTest} de Auth, Usuarios, Académico y
Reservas ejecutan esta comprobación en cada \texttt{mvn verify}; Auth incluye
además \texttt{SchemaContractVerifierNegativeTest} para mutaciones negativas
del verificador.

El verificador contrasta campo por campo los DTO de request y response
publicados: tipos básicos, UUID, fechas \texttt{date}/\texttt{date-time},
números, arrays y colecciones, DTO anidados mediante \texttt{\$ref}, enums y
obligatoriedad contextual según uso como request/response y evidencia de Bean
Validation o serialización de Jackson. Para paginación exige un objeto con las
nueve propiedades de \texttt{Page} o todos los componentes del record
\texttt{PaginaResponse}, derivados por reflexión, y comprueba que el contenido
sea un array con referencia al tipo de elemento \texttt{T}. El Gateway no tiene
DTO propios equivalentes en esta verificación; su contrato se deriva de los
servicios y su completitud se valida por rutas, alias y catálogo ejecutable.
Esta validación no pretende demostrar reglas de negocio, ejemplos,
descripciones ni toda posible restricción semántica de OpenAPI. La única fuente canónica de \texttt{SchemaContractVerifier} se encuentra en
\path{tests/openapi-runtime/src/main/java/ec/edu/scli/contracts/SchemaContractVerifier.java}.
Los cuatro servicios incorporan esa carpeta como fuente de pruebas con
\texttt{build-helper-maven-plugin:add-test-source} en
\texttt{generate-test-sources}, sin copias ni symlinks. También se eliminan las
cuatro copias locales idénticas de \texttt{RuntimeContractVerifier} para evitar
clases duplicadas al compilar la carpeta compartida.
El job \texttt{validate-openapi-contracts} ejecuta las pruebas Python del
generador además de la validación estructural y condiciona
\texttt{publish-release} mediante \texttt{needs}, conservando sus dependencias
anteriores. El trigger de push incluye tags \texttt{v*} y el job OpenAPI no
tiene una condición que los excluya.

La validación previa acreditó cuatro mutaciones negativas, todas revertidas:
(1) un objeto paginado con propiedades falsas; (2) eliminar
\texttt{totalPaginas} de \texttt{PaginaResponse}; (3) cambiar el default de
\texttt{pagina} por un valor incorrecto; y (4) añadir un
\texttt{RouterFunction} con una ruta arbitraria que omitía el catálogo.
Las tres primeras hicieron fallar el verificador semántico; la cuarta hizo
fallar la garantía arquitectónica del Gateway.

OpenAPI describe correctamente la superficie, los parámetros de consulta
relevantes, los cuerpos request/response y los schemas DTO campo por campo solo
a partir de esta corrección. La estructura general del contrato se valida con
las comprobaciones OpenAPI y de consistencia del snapshot; la afirmación
independiente adicional cubre los contratos HTTP y sus DTO publicados, sin
extenderse a semántica funcional ni reglas de negocio.

La validación local de este cierre, en Windows y Java 21, ejecutó 41 pruebas
correctas del Gateway y una prueba \texttt{OpenApiSchemaContractTest} correcta
por cada servicio, compilando la fuente compartida. Las cinco pruebas Python
del generador pasaron y el validador estructural terminó sin errores.
Dos mutaciones temporales de producción, una construcción alternativa con
consulta al catálogo y otra sin helper, hicieron fallar el guard con los
nombres \texttt{bypassDocenteTemporal} y \texttt{sinHelperTemporal}.
Ambas se retiraron antes de la ejecución final satisfactoria.
